% !TeX root = ../教材题库_全部条目_转写工作稿.tex
% Source: XJ-B1-C06-RV-P266-Y1140
% 第6章 / 章末复习 复习题六 / 教材 p.266、267 / PDF 第273、274页
% OCR未校对：公式与文字均可能有误；未完成数学TeX转写；禁止直接发给学生
% 原题及相邻原解截图（核对用）：../题目原图/XJ-B1-C06-RV-P266-Y1140-01.webp, ../题目原图/XJ-B1-C06-RV-P266-Y1140-02.webp
% 本文件只在转写完成后作为题干使用；OCR 条目尚不是可靠题干。
\begin{bankitem}{复习题六 第12题}{第6章 / 章末复习 复习题六 / 教材 p.266、267 / PDF 第273、274页}{XJ-B1-C06-RV-P266-Y1140}{OCR未校对：公式与文字均可能有误；未完成数学TeX转写；禁止直接发给学生}
\begin{ocrdraft}
[教材p.266 / PDF第273页]
\par
12．立德中学从参加高三模拟考试的学生中随机抽取60名学生，将其数学成
\par
绩（均为整数）分成六段[90，100），[100，110)，
\par
040，150]后得到如下部
\par
分频率分布直方图．
\par
组距
\par
0．025
\par
0．020
\par
0，015
\par
0．010
\par
0．005
\par
0
\par
90100110120130140150《分娄
\par
（第12题）
\par
[教材p.267 / PDF第274页]
\par
观察图形，回答下列问题：
\par
（D求分数在[120，130）内的频率，并补全这个频率分布直方图；
\par
（2）估计本次考试成绩的平均数、众数、中位数．
\end{ocrdraft}
\end{bankitem}
